\section{Introdução}
O desenvolvimento de software moderno demanda níveis elevados de qualidade, confiabilidade e agilidade para atender a cenários de rápida evolução tecnológica. Em projetos de médio a grande porte, a complexidade do código-fonte cresce exponencialmente, tornando imprescindível a utilização de práticas de verificação contínua por meio de testes sistemáticos. Neste contexto, os testes unitários representam a primeira linha de defesa contra defeitos de implementação, pois validam o comportamento de métodos ou funções isoladas, detectando falhas ainda na fase inicial de codificação e evitando que erros se propaguem para estágios posteriores, como integraçao e homologação.

Todavia, a elaboração manual de uma suíte abrangente de testes unitários é tarefa trabalhosa, repetitiva e suscetível a omissões humanas. Ferramentas tradicionais de geração automática de testes, como o EvoSuite [3], empregam técnicas baseadas em análise estática, geração evolutiva ou fuzzing para produzir casos de teste em Java, mas apresentam limitações quando o método-alvo possui lógica complexa, com múltiplos laços aninhados, ramificações condicionais profundas ou uso de recursos avançados da linguagem, como reflexão e streams. Além disso, atender a critérios avançados de cobertura - tais como cobertura de todos os caminhos (all-paths), de todas as definições-uso (all-defs-uses) e de todos os usos (all-uses) - continua sendo um desafio para abordagens convencionais.

Nos últimos anos, avanços em inteligência artificial (IA), especialmente no campo dos Modelos de Linguagem de Grande Porte (LLMs), demonstraram capacidade semântica considerável para compreensão e geração de código em diversas linguagens de programaçãoo. Ferramentas como GitHub Copilot (baseado em OpenAI Codex) [1], Code Llama e variantes especializadas de GPT-4 foram capazes de sugerir implementações, corrigir trechos de código e até gerar testes unitários com base em descrições em linguagem natural ou diretamente a partir do código-fonte. Pesquisas recentes apontam para o potencial dos LLMs na criação automática de testes com qualidade [2, 8, 5, 13]. No entanto, ainda há lacunas quanto à garantia de que esses testes gerem cobertura completa segundo critérios formais de análise de código.

\subsection{Justificativa}
Diante desses avanços, este trabalho concentra-se na investigação da eficácia dos LLMs na geração automática de testes unitários para métodos Java, com ênfase em alcançar critérios avançados de cobertura de código. A relevância deste estudo reside na possibilidade de automatizar uma etapa crítica do desenvolvimento de software, reduzindo o esforço manual e aumentando a confiabilidade dos testes. Além disso, compreender as limitações e potencialidades dos LLMs nesse contexto pode orientar futuras pesquisas e aprimoramentos em ferramentas de geração automática de testes, contribuindo para a evolução das práticas de desenvolvimento ágil e DevOps.

\subsection{Objetivos}
Sendo assim, este trabalho tem como objetivo a investigar a capacidade de LLMs na geração automática de testes unitários para métodos Java, avaliando a cobertura de código alcançada e comparando os resultados tradicionais de geração de testes, como o EvoSuite.